\section{Geography}

\subsection{Geographic Attributes}

\subsection{Cities}

\subsection{Rating Guide}
Considering these geographical factors allows for a
comprehensive evaluation of a country's unique
characteristics and challenges, providing insights into
its potential for economic development, environmental
sustainability, and geopolitical significance.
\begin{enumerate}
  \item \textbf{Location and Borders}: The country's position
  on the globe and its borders with neighboring nations, which
  can influence trade, security, and geopolitical dynamics.
  
  \item \textbf{Land Area and Terrain}: The total land area
  and the diversity of terrain, including mountains, plains,
  forests, and water bodies, impacting agriculture, biodiversity,
  and infrastructure development.
  
  \item \textbf{Natural Resources}: The availability of natural
  resources, such as minerals, oil, water, and fertile soil,
  affecting economic development and resource management.
  
  \item \textbf{Coastline and Ports}: The length of the coastline
  and the presence of ports, crucial for maritime trade, fishing,
  and coastal development.
  
  \item \textbf{Natural Hazards}: The occurrence of natural
  disasters like earthquakes, tsunamis, hurricanes, and floods,
  affecting infrastructure resilience and disaster preparedness.
  
  \item \textbf{Environmental Factors}: The level of pollution, air
  quality, deforestation, and conservation efforts, influencing the
  overall environmental health of the country.
  
  \item \textbf{Geopolitical Significance}: The country's strategic
  importance in regional and international politics, trade routes,
  and alliances.
  
  \item \textbf{Geographical Diversity}: The range of landscapes
  and ecosystems within the country, contributing to tourism,
  biodiversity, and cultural heritage.
  
  \item \textbf{Disputed Territories}: Any contested or disputed
  areas within the country's borders, influencing stability and
  international relations.
\end{enumerate}


